\documentclass[10pt,a4paper]{amsart}
\usepackage[margin=19mm]{geometry}
\usepackage{amsmath,amssymb,mathtools}
\usepackage[scr=rsfs,cal=euler]{mathalfa}
\usepackage{xfrac}
\usepackage[unicode,hidelinks,bookmarks=false]{hyperref}
\newcommand{\ie}{i.e.\ }
\newcommand{\cf}{cf.\ }
\newcommand{\resp}{respectively\ }
\newcommand{\Subsection}{\subsection}
\providecommand{\nobreakdash}{\nobreak\hbox{-}}
\setlength{\emergencystretch}{2em}
\multlinegap0pt
\usepackage{amssymb}
\usepackage{enumerate}
\usepackage{xcolor}
\newtheorem{theorem}{Theorem}
\newtheorem{remark}{Remark}
\newtheorem{lemma}{Lemma}
\newcommand{\calB}{\mathcal{B}}
\newcommand{\calH}{\mathcal{H}}
\newcommand{\cut}{\backslash\backslash}
\title{Some fast algorithms for curves in surfaces}
\author{Marc Lackenby}
\date{Source: arXiv:2401.16056v2 (CC BY 4.0)}
\begin{document}
\maketitle

\noindent\emph{Source and licence.} Some fast algorithms for curves in surfaces. Version arXiv:2401.16056v2 (CC BY 4.0), distributed by arXiv under the Creative Commons Attribution 4.0 International licence, http://creativecommons.org/licenses/by/4.0/, as recorded in the arXiv licence metadata for this exact version and shown on the abstract page https://arxiv.org/abs/2401.16056v2. Full licence notice: \texttt{licenses/arxiv-2401.16056-CC-BY-4.0.txt}.\par\smallskip

\noindent\textit{Editorial scope.} Counted unit: the article's theorem Thm:Normalise1Manifold (source0.tex lines 487--495). The article's complete original proof of it is delivered with it (source0.tex lines 1134--1161, one \texttt{proof} environment of 28 lines, which the article places away from the statement). No mathematics is written, repaired or re-derived: the statement and the proof are the article's own lines, in the article's own notation, and no retained line is substituted or re-flowed.  Retained uncounted as the source-local context the proof needs: lines 736--739; lines 757--760; lines 762--775; lines 983--986; lines 998--1009; lines 1037--1040; lines 1051--1054; lines 1077--1082; lines 1112--1115; lines 1122--1125.  One declared adaptation: exactly 1 ancillary strings inside the retained spans are omitted (image-inclusion commands and, where the source repeats a label, the repeated label definitions) and no other line differs from the source; the figure environments, with their placement, captions and labels, are kept, and the package records the omitted strings in fidelity receipts bound to the affected bodies.  No retained line contains a citation, so the wrapper carries no bibliography and no bibliography entry.  The retained lines contain the article's own diagram code, which the wrapper typesets with the standard tikz packages; no image file is delivered and the package declares no asset dependency.  Editorial formatting only, in addition: an AMS \texttt{amsart} preamble that restates the article's own declarations and macros used by the retained lines, namely the environments \texttt{theorem}, \texttt{remark}, \texttt{lemma} on separate counters, together with the standard packages those lines need.  The retained environments print in this document as Theorem 1, Remark 1, Lemma 1 in order of appearance, whereas the article prints the same items with the numbering of its own full document.  The complete native source of the article as distributed in the arXiv e-print for this exact version is bundled unchanged as \texttt{sources/153667/source0.tex}.
\begin{theorem}[{\sc Cut along standard 1-manifold}]
\label{Thm:CutAlong1Manifold}
There is an algorithm that takes, as its input, a handle structure $\calH$ for a compact surface $S$ and a standard 1-manifold $C$ properly embedded in $S$, and outputs the reduced handle structure on the relevant components of $S \cut C$. The running time is at most a polynomial function of $||\calH||$, $\mathrm{simp}(C)$, and $\log (w(C) + |\partial C|)$.
\end{theorem}

\begin{remark} 
\label{Rem:ReducedHSEmbeds}
It is worth stating that the reduced handle structure $\calH'$ on $S \cut C$ is provided not just as an abstract handle structure but also with information about how it embeds in $\calH$. More precisely, for each handle of $\calH'$ that did not come from a component of the parallelity bundle $\calB$, we have an identification between it and the handle of $S \cut C$ that it came from. We can also determine the location of the handles arising from components of $\calB$, in the following sense.
\end{remark}

\begin{theorem}[{\sc Location of parallelity bundle component}]
\label{Thm:LocationParBundle}
There is an algorithm that takes, as its input, the following:
\begin{enumerate}
\item a handle structure $\calH$ for a compact surface $S$;
\item a standard 1-manifold $C$ properly embedded in $S$;
\item a point $p$ in the intersection between $S - C$ and a 1-cell of the associated cell structure on $S$;
\end{enumerate} and provides the following:
\begin{enumerate}
\item the vector(s) in $S$ for the component(s) of $\partial_hB$, where $B$ is the component of the parallelity bundle for $S \cut C$ containing $p$ (or  the zero vector if $p$ does not lie in the parallelity bundle);
\item for each such component $\alpha$ of $\partial_h B$, the location of $\partial \alpha - \partial S$.
\end{enumerate}
The running time is at most a polynomial function of $||\calH||$, $\mathrm{simp}(C)$, and $\log (w(C) + |\partial C|)$.
\end{theorem}

\begin{lemma}
\label{Lem:NotSelfReturning}
Let $\calH$ be a handle structure for a compact surface $S$. Let $C$ be a standard 1-manifold properly embedded in $S$. Let $D$ be a simplifying disc for $C$. Suppose that $D$ is not self-returning. Then after performing at most $18||\calH|| + 4\mathrm{simp}(C)$ isotopies across maximal enlargements of isotopy regions containing $D$, the resulting 1-manifold $C''$ satisfies $\mathrm{simp}(C'') < \mathrm{simp}(C)$.
\end{lemma}

\begin{lemma}[{\sc Isotopy across maximal enlargement}]
\label{Lem:IsotopyAcrossMaximalEnlargement}
There is an algorithm that takes as its input the following:
\begin{enumerate}
\item a handle structure $\calH$ for a compact surface $S$;
\item a standard 1-manifold $C$ properly embedded in $S$;
\item an arc of $C \cap \calH^0$ transversely oriented into a simplifying disc $D$;
\end{enumerate}
and outputs the standard 1-manifold $C'$ that is obtained by performing the isotopy across the maximal enlargement of the isotopy region associated with $D$,
and possibly removing some boundary-parallel arc components. It also outputs the number of arcs of intersection between $C$ and this maximal enlargement.
The running time is bounded above by a polynomial function of $||\calH||$, $\mathrm{simp}(C)$, and $\log (w(C) + |\partial C|)$.
\end{lemma}

\begin{lemma}
\label{Lem:SelfReturningHoriz}
Let $\calH$ be a handle structure for a compact surface $S$. Let $C$ be a standard 1-manifold properly embedded in $S$. Let $D$ be a simplifying disc for $C$, and let $R$ be the associated isotopy region. Suppose that the horizontal boundary of $R$ contains a standard arc in $\calH^0$ that is parallel to $D \cap C$. Then after performing at most $36||\calH|| + 8 \mathrm{simp}(C)$ isotopies across maximal enlargements of isotopy regions containing $D$, the resulting 1-manifold $C''$ satisfies ${\mathrm{simp}}(C'') < {\mathrm{simp}}(C)$.
\end{lemma}

\begin{lemma}
\label{Lem:SelfReturningNotReeb}
Let $\calH$ be a handle structure for a compact surface $S$. Let $C$ be a standard 1-manifold properly embedded in $S$. Let $D$ be a simplifying disc for $C$, and let $R$ be the associated isotopy region. Suppose $R$ is self-returning and that no arc of $\partial_hR$ is normally parallel to $D \cap C$. Let $D_+$ be the maximal enlargement of $R$. Suppose that $|D_+ \cap C|$ is strictly less than the number of arcs of $C \cap \calH_0$ normally parallel to $D \cap C$, including $D \cap C$. Let $C'$ be the 1-manifold obtained from $C$ by the isotopy across this maximal enlargement. Then $C' \cap \calH^0$ has a component that is parallel to $D \cap C$ and that bounds a simplifying disc for $C'$ which is not self-returning.
\end{lemma}

\begin{figure}[h]
\centering

\caption{A maximal Reeb region $A$ is shaded, and the associated Reeb isotopy is shown. The left and right sides of these figures are identified, but possibly via a reflection, in which case $A$ would be a M\"obius band}
\label{Fig:ReebRegion}
\end{figure}

\begin{lemma}
\label{Lem:ReebConclusion}
Let $\calH$ be a handle structure for a compact surface $S$. Let $C$ be a standard 1-manifold properly embedded in $S$. Let $D$ be a simplifying disc for $C$, and let $R$ be the associated isotopy region. Suppose $R$ is self-returning and that no arc of $\partial_hR$ is normally parallel to $D \cap C$. Let $D_+$ be the maximal enlargement of $R$. Suppose that $|D_+ \cap C|$ is not strictly less than the number of arcs of $C \cap \calH_0$ normally parallel to $D \cap C$, including $D \cap C$. Then $D_+$ lies in a Reeb or semi-Reeb region for $C$.
\end{lemma}

\begin{lemma}
\label{Lem:ReebIsotopy}
Let $\calH$ be a handle structure for a compact surface $S$. Let $C$ be a standard 1-manifold properly embedded in $S$. Let $D$ be a simplifying disc for $C$, and suppose that $D$ lies in a maximal Reeb or semi-Reeb region $A$. Let $C'$ be the 1-manifold obtained by the Reeb or semi-Reeb isotopy, and let $D'$ be the resulting simplifying disc for $C'$. Suppose that $D'$ is self-returning. Then $D'$ satisfies the hypotheses of Lemma \ref{Lem:SelfReturningNotReeb}.
\end{lemma}

\begin{theorem}[{\sc Normalise a standard 1-manifold in a handle structure}]
\label{Thm:Normalise1Manifold}
There is an algorithm that takes as its input the following:
\begin{enumerate}
\item a handle structure $\calH$ for a compact surface $S$;
\item a standard 1-manifold $C$ properly embedded in $S$;
\end{enumerate}
and outputs a normal 1-manifold $\overline{C}$ that is obtained from $C$ by an isotopy plus possibly discarding some curves that bound discs and boundary-parallel arcs. The running time for the algorithm is bounded above by a polynomial function of $||\calH||$, $\mathrm{simp}(C)$ and $\log (w(C) + |\partial C|)$.
\end{theorem}

\begin{proof}[Proof of Theorem \ref{Thm:Normalise1Manifold}]
We perform the following operations in a loop. First, we search for the existence of a maximal Reeb or semi-Reeb region. We will explain how to do this below. If there is such a Reeb or semi-Reeb region, we perform this Reeb or semi-Reeb isotopy, creating a 1-manifold $C'$. Denote the simplifying disc for $C'$ in the Reeb or semi-Reeb region by $D'$. 
If there was no Reeb or semi-Reeb region, then we let $C' = C$ and we pick any simplifying disc $D'$ for $C'$. 

We then perform the isotopy across the maximal enlargement of the isotopy region associated with $D'$. If this reduces the simplification number, then we return to the start of the loop. On the other hand, if the simplification number is unchanged, then we know that the new 1-manifold $C''$ has an arc of $C'' \cap \calH^0$ that runs parallel to $D' \cap C'$. An outermost such arc bounds a simplifying disc. We repeat the procedure in this paragraph with this disc. We continue until the simplification number decreases, at which point we return to the start of the loop. This continues until there are no more simplifying discs.

We now explain why this procedure works. If there is no Reeb or semi-Reeb region, then, by Lemma \ref{Lem:ReebConclusion}, each simplifying disc $D'$ for $C'$ satisfies at least one of the following: 
\begin{enumerate}
\item $D'$ is not self-returning;
\item the horizontal boundary of its isotopy region contains a standard arc that runs parallel to $D' \cap C'$;
\item the horizontal boundary of the isotopy region $R'$ does not contain a standard arc that runs parallel to $D' \cap C'$, and the number of arcs parallel to $D' \cap C'$ (including $D' \cap C'$) is greater than the number of arcs of $C'$ in the maximal enlargement of $R'$.
\end{enumerate}
In the first case, Lemma \ref{Lem:NotSelfReturning} gives that after performing at most $18||\calH|| + 4\mathrm{simp}(C)$ isotopies across maximal enlargements, the simplification number goes down. In the second case, a similar conclusion holds by Lemma \ref{Lem:SelfReturningHoriz}. In the third case, Lemma \ref{Lem:SelfReturningNotReeb} implies that the 1-manifold obtained by isotoping $C'$ across the maximal enlargement of $R'$ has a simplifying disc that is not self-returning. Hence, after at most $18||\calH|| + 4\mathrm{simp}(C)$ further isotopies, the simplification number goes does.
On other hand, if there is a Reeb or semi-Reeb region for $C$, then there is a maximal one $A$, and we let $C'$ be the 1-manifold obtained by the Reeb or semi-Reeb isotopy and we let $D'$ be the simplifying disc for $C'$ contained in $A$. Lemma \ref{Lem:ReebIsotopy} and Lemma \ref{Lem:SelfReturningNotReeb} imply that either $D'$ is not self-returning or the 1-manifold obtained by the isotopy across the maximal enlargement of $D'$ has such a simplifying disc. Hence, after at most $18||\calH|| + 4\mathrm{simp}(C)$ further isotopies, the simplification number decreases.

We now describe how to find a maximal Reeb region if one exists. We can readily find all the interior-simplifying discs for the 1-manifold $C$, by considering each parallelism class of arcs of $C \cap \calH^0$. We consider each interior-simplifying disc in turn. So let us focus on one such disc $D$ and its isotopy region $R$. We first determine whether $R$ is self-returning as follows. By applying {\sc Cut along standard 1-manifold} (Theorem \ref{Thm:CutAlong1Manifold}) to $S \cut C$ and Remark \ref{Rem:ReducedHSEmbeds}, we can  determine the component $W$ of $\calH^0 \cut C$ containing $\partial_v R \cut D$. (See Figure \ref{Fig:ReebRegion}.) Then $R$ is self-returning if and only if the following all hold:
\begin{enumerate}
\item $W$ lies in the interior of $S$;
\item the intersection between $\partial W$ and $\calH^0 \cap C$ consists of three arcs $\beta_-$, $\beta_0$ and $\beta_+$;
\item one of these arcs $\beta_0$ is normally parallel to $D \cap C$ and all arcs of $C \cap \calH^0$ between $\beta_0$ and $D \cap C$ are normally parallel;
\item the remaining two arcs $\beta_-$ and $\beta_+$ of $\calH^0 \cap C \cap \partial W$ are normally parallel to each other.
\end{enumerate}
These conditions are easily checked. Suppose that $R$ is self-returning. Then we can easily check whether the horizontal boundary of $R$ contains no arc of $\calH^0 \cap C$ that is normally parallel to $D \cap C$ using {\sc Location of parallelity bundle component} (Theorem \ref{Thm:LocationParBundle}). Using {\sc Isotopy across maximal enlargement} (Lemma \ref{Lem:IsotopyAcrossMaximalEnlargement}), we can determine the number $p$ of arcs of intersection between $C$ and the maximal enlargement of $R$. For $R$ to be part of a Reeb region, this number needs to be equal to the number of components of $C \cap \calH^0$ normally parallel to $C \cap D$, including $C \cap D$ itself. 

Thus, using the above procedure, we can determine whether $R$ is part of a Reeb region. Suppose it is. We then find the maximal Reeb region containing $R$, as follows. Let $\beta$ be the component of $\calH^0 \cap \calH^1$ incident to $D$. As in the proof of Lemma \ref{Lem:IsotopyAcrossMaximalEnlargement}, we can determine whether $R$ intersects $\beta - D$. If it does not, then we cut $C$ along $\beta$ and count the number of components of $C - \beta$ normally parallel to each component of $\partial_h R$. Let $q_1$ and $q_2$ be these number of arcs. Let $q$ be the largest integer multiple of $p$ that is at most $\min \{ q_1, q_2 \}$. Then the union of this many copies of $\partial_h R$, together with all arcs of $C \cap \calH^0$ parallel to $C \cap D$, forms the arcs of intersection between $C$ and the maximal Reeb region $A$. Removing all these arcs and replacing them with this many copies of $C \cap D$ achieves the Reeb isotopy. 

There is a similar procedure for finding a maximal semi-Reeb region if one exists and for performing the associated semi-Reeb isotopy.
\end{proof}

\end{document}
